% !TEX root = ../../main_without_quantization.tex
\subsection{Algorithms for Gradient-Based Optimization}
\label{sec:gradient_based_optimization}
With the loss function established, the remaining question is how the resulting error signal propagates back through the network to update the parameters.

Training neural networks is usually a problem of finding a parameter vector $\theta \in \mathbb{R}^{d}$ that minimizes the scalar objective $J(\theta)$ where $J$ is the empirical risk defined by the loss function and the training data. Formally, we want to solve \eqref{eq:optimization_problem}:
\begin{equation}
    \theta^* = \arg\min_{\theta} J(\theta).
    \label{eq:optimization_problem}
\end{equation}

The parameter space in most contemporary neural networks is high-dimensional and the objective is largely non-convex so a solution in closed form for optimum parameters rarely exists. Training is hence iterative: the parameters are updated according to directions suggested by the gradient of the loss \cite{Bottou2018}.

The gradient $\nabla J(\theta)$ indicates the direction of steepest local increase of the objective, and optimizers will use it in reverse to adjust the parameters to  yield a lower loss. This concept forms the basis of the gradient descent algorithm as well as of all the adaptive optimization schemes currently used in training neural networks.

\subsubsection{Gradient Descent}

Gradient descent updates parameters in the negative-gradient direction. For iteration $t$ and learning rate $\eta>0$, the update is formulated in Eq.\eqref{eq:gradient_descent}:
\begin{equation}
    \label{eq:gradient_descent}
    \theta_{t+1} = \theta_t - \eta \nabla J(\theta_t).
\end{equation}

The learning rate controls the step size. If it is too small, optimization may progress slowly; if it is too large, the updates can overshoot useful regions of the objective surface. In large scale learning, the exact gradient over the full dataset is often replaced by a stochastic estimate computed on a mini-batch. This gives stochastic gradient descent, which is cheaper per update and is widely used for neural-network training \cite{Bottou2018}.

\subsubsection{Adam}

Adam is an adaptive gradient-based optimizer that combines momentum with per-parameter learning-rate adaptation \cite{kingma2015}. Instead of using only the current gradient, Adam keeps exponential moving averages of the first and second moments of the gradient, which are updated as shown in Eq.\eqref{eq:adam_moments}:
\begin {equation}
    \label{eq:adam_moments}
    m_t = \beta_1 m_{t-1} + (1-\beta_1)g_t,
    \qquad
    v_t = \beta_2 v_{t-1} + (1-\beta_2)g_t^2,
\end{equation}

%\begin{align}
%    m_t &= \beta_1 m_{t-1} + (1-\beta_1)g_t,\\
%    v_t &= \beta_2 v_{t-1} + (1-\beta_2)g_t^2,
%\end{align}
where $g_t=\nabla J(\theta_t)$ is the gradient estimate at iteration $t$, $m_t$ tracks the average direction of recent gradients, and $v_t$ tracks their squared magnitude.

Because both moving averages are initialized at zero, Adam applies bias correction computed as follows Eq.\eqref{eq:adam_bias_correction}:
\begin{equation}
    \label{eq:adam_bias_correction}
    \hat{m}_t = \frac{m_t}{1-\beta_1^t},
    \qquad
    \hat{v}_t = \frac{v_t}{1-\beta_2^t}.
\end{equation}
The parameter update is then given by \eqref{eq:adam_update}:
\begin{equation}
    \label{eq:adam_update}
    \theta_{t+1}
    =
    \theta_t
    -
    \eta \frac{\hat{m}_t}{\sqrt{\hat{v}_t}+\epsilon}.
\end{equation}

Adam is often utilized due to its robustness in relation to scaling differences among gradients for different parameters and typically requires less fine-tuning of learning rates as opposed to simple gradient descent. Subsequent versions of Adam include AdamW, which removes the weight decay from the adaptive learning component and is applicable to a variety of deep learning applications \cite{Loshchilov2019}.



Both gradient descent and Adam describe how the parameters should be updated, but they both require that the gradients be available. These computations can become expensive in deep, multi-layered networks and thus require an alternative method, termed backpropagation.